\documentclass[11pt, a4paper]{article}

% --- UNIVERSAL PREAMBLE BLOCK ---
\usepackage[a4paper, top=2.5cm, bottom=2.5cm, left=2.54cm, right=2.54cm]{geometry}
\usepackage{fontspec}

% Allow long URLs in the bibliography to break seamlessly
\usepackage{xurl}      

\usepackage[english, main=hindi]{babel}

% Primary font configuration.
\directlua{luaotfload.add_fallback("englishfallback", {
    "Noto Sans:mode=node;",
    "Noto Sans Bengali:mode=node;"
})}

\babelfont{rm}[
    Script=Devanagari,
    Renderer=Node,
    RawFeature={fallback=englishfallback},
    AutoFakeSlant=0.2,
    ItalicFont=Noto Sans Devanagari,
    ItalicFeatures={FakeSlant=0.2},
    BoldItalicFont=Noto Sans Devanagari,
    BoldItalicFeatures={FakeSlant=0.2}
]{Noto Sans Devanagari}

% --- ADDITIONAL REQUIRED PACKAGES ---
\usepackage{amsmath}
\usepackage{amssymb}
\usepackage{booktabs}
\usepackage{tabularx}
\usepackage{enumitem}
\usepackage{setspace}
\usepackage{indentfirst} 
\usepackage{titlesec}
\usepackage{caption}

% Caption styling compliant with APA 7th Standards
\captionsetup[table]{
    position=above,
    justification=raggedright,
    singlelinecheck=false,
    labelfont=bf,
    textfont=it,
    skip=6pt
}

% APA 7th Edition Spacing and Paragraph Indentation
\setstretch{1.4}
\setlength{\parindent}{1.27cm}
\setlist[itemize]{label=--, leftmargin=1.8cm}
\setlist[enumerate]{leftmargin=1.8cm}

% APA 7th Edition Heading Formats
\setcounter{secnumdepth}{0}
\titleformat{\section}[block]{\normalfont\normalsize\bfseries\centering}{}{0pt}{}
\titleformat{\subsection}[block]{\normalfont\normalsize\bfseries}{}{0pt}{}
\titleformat{\subsubsection}[runin]{\normalfont\normalsize\bfseries\itshape}{}{0pt}{}[.]

% APA Citation Formatting in References
\usepackage[round, authoryear]{natbib}
\setcitestyle{aysep={,}}
\setlength{\bibhang}{1.27cm}

\begin{document}

% --- TITLE PAGE (APA 7th Student Standard) ---
\begin{titlepage}
    \centering
    \vspace*{1.5in}
    
    {\Large \textbf{भारत में भ्रष्टाचार निवारण हेतु इतिहास में उठाए गए प्रभावी कदम:}} \\[0.5em]
    {\large \textbf{एक खेल-सैद्धांतिक, मनोवैज्ञानिक एवं अनुभवजन्य विश्लेषण}} \\[2.5\baselineskip]
    
    {\normalsize \textbf{शोधार्थी / पंजीयन क्रमांक:} REG-2026-XXXX} \\[0.4em]
    {\normalsize शासकीय उत्कृष्ट उच्चतर माध्यमिक विद्यालय, जबलपुर} \\[0.4em]
    {\normalsize कक्षा १२ (द्वादश) -- विज्ञान एवं सामाजिक विज्ञान संकाय} \\[2.5\baselineskip]
    
    {\normalsize प्रतिभा सम्मान योजना (PSY) -- राष्ट्रीय शोध प्रतियोगिता (कक्षा ६ से १२ वर्ग)} \\[0.4em]
    {\normalsize निर्धारित शोध विषय क्रमांक: १} \\[0.4em]
    {\normalsize शैक्षणिक सत्र: २०२६--२०२७}
    
    \vfill
    \noindent \rule{0.6\textwidth}{0.4pt} \\[0.5em]
    {\footnotesize \textit{प्रस्तुत शोध-प्रबंध प्रतिभा सम्मान योजना (PSY) के मूल्यांकन मानकों तथा APA ७वें संस्करण के औपचारिक शोध प्रारूप के पूर्ण अनुपालन में प्रस्तुत है।}}
    \vspace*{0.5in}
\end{titlepage}

% --- ABSTRACT PAGE ---
\newpage
\begin{center}
    \textbf{सारांश (Abstract)}
\end{center}

\noindent पारंपरिक प्रशासनिक एवं सामाजिक विमर्श में भ्रष्टाचार को ऐतिहासिक रूप से केवल व्यक्तिगत नैतिकता, देशभक्ति की कमी और चारित्रिक पतन के रूप में चित्रित किया जाता रहा है। इसके प्रत्युत्तर में नीतिगत स्तर पर उपदेशात्मक आचार संहिताओं, सत्यनिष्ठा प्रतिज्ञाओं और नैतिक व्याख्यानों पर अत्यधिक बल दिया गया। प्रस्तुत शोध-प्रबंध मूल्य-तटस्थ संस्थागत अर्थशास्त्र और अनुभवजन्य साक्ष्यों के आधार पर यह विश्लेषित करता है कि शुद्ध नैतिक दृष्टिकोण प्रशासनिक धरातल पर प्रणालीगत सुधार लाने में संरचनात्मक रूप से क्यों सीमित रहा है। अध्ययन में रॉबर्ट क्लिटगार्ड के किराया-वसूली प्रतिमान और जॉन नैश के खेल सिद्धांत (Game Theory) की सुगम्य व्याख्या करते हुए यह सिद्ध किया गया है कि अत्यधिक प्रक्रियात्मक विलंब और असीमित प्रशासनिक विवेकाधिकार के वातावरण में नागरिक और लोकसेवक दोनों के लिए रिश्वत का लेन-देन एक गणितीय रूप से प्रभावी 'उप-इष्टतम नैश संतुलन' बन जाता है। इसके समानांतर, अत्यधिक प्रतिस्पर्धी चयन प्रक्रियाओं से उत्पन्न डूबी हुई लागत का भ्रम (Sunk Cost Fallacy), अधिकार-बोध (Psychological Entitlement) और नैतिक अलगाव (Moral Disengagement) जैसे मनोवैज्ञानिक प्रेरकों का विशद परीक्षण किया गया है। अध्ययन में कौटिल्य के अर्थशास्त्र, संथानम समिति (१९६४), भ्रष्टाचार निवारण अधिनियम (१९८८/२०१८), सूचना का अधिकार (RTI, २००५) और प्रत्यक्ष लाभ अंतरण (DBT) के ऐतिहासिक प्रभाव का मूल्यांकन किया गया है। इन संरचनात्मक हस्तक्षेपों का स्थानीय स्तर पर सत्यापन करने हेतु एक प्राथमिक विद्यालयी पारिस्थितिकी तंत्र ($N=60$) में संरचित क्षेत्रीय सर्वेक्षण आयोजित किया गया। अंततः, प्रशासनिक विवेकाधिकार विलोपन, स्वतः-स्वीकृति प्रणाली (Deemed Approval) और वित्तीय मुखबिरी पुरस्कारों पर आधारित डेटा-समर्थित नीतिगत सुधार प्रस्तुत किए गए हैं।

\vspace{1\baselineskip}
\noindent \textit{मूल शब्द (Keywords):} संस्थागत अर्थशास्त्र, खेल सिद्धांत (Game Theory), नैश संतुलन (Nash Equilibrium), किराया वसूली (Rent-Seeking), डूबी हुई लागत (Sunk Cost), प्रत्यक्ष लाभ अंतरण (DBT)।

% --- MAIN BODY ---
\newpage
\begin{center}
    \textbf{भारत में भ्रष्टाचार निवारण हेतु इतिहास में उठाए गए प्रभावी कदम: एक खेल-सैद्धांतिक, मनोवैज्ञानिक एवं अनुभवजन्य विश्लेषण}
\end{center}

लोक प्रशासन का प्राथमिक उद्देश्य जनहित का निष्पक्ष संपादन और राज्य के संसाधनों का न्यायसंगत वितरण सुनिश्चित करना है। तथापि, विकासशील राष्ट्रों के प्रशासनिक इतिहास में भ्रष्टाचार एक चिरस्थायी संरचनात्मक व्याधि के रूप में विद्यमान रहा है। पारंपरिक सामाजिक विमर्श में इस समस्या को प्रायः व्यक्तिगत चारित्रिक दोष मानकर देखा जाता रहा है, जिसके कारण इसके समाधान भी उपदेशात्मक सीमाओं तक सिमट कर रह गए। प्रस्तुत शोध-पत्र भ्रष्टाचार को नैतिक पतन के रूप में देखने के बजाय संस्थागत विसंगतियों, रणनीतिक लाभ-हानि निर्णयों और संज्ञानात्मक पूर्वाग्रहों के वैज्ञानिक संगम के रूप में विश्लेषित करता है।

\section{सैद्धांतिक एवं वैचारिक पृष्ठभूमि: नैतिक बनाम संरचनात्मक दृष्टिकोण}

\subsection{पूर्व पक्ष: व्यक्तिगत चरित्र एवं नैतिक-आदर्शवादी दृष्टिकोण का सैद्धांतिक आधार}

प्रशासनिक दर्शन और राजनीतिक आचारशास्त्र के प्रारंभिक साहित्य में भ्रष्टाचार को मूलतः एक 'व्यक्तिगत नैतिक विचलन' (Individual Moral Transgression) तथा सार्वजनिक पद के न्यास (Trust) के प्रति विश्वासघात के रूप में परिभाषित किया गया \citep{dobel1978, philp1997, thompson1993}। इस शास्त्रीय परंपरा के अनुसार, राज्य एक नैतिक संस्था है और लोकसेवक उसका संवाहक; अतः पद का दुरुपयोग लोकसेवक के व्यक्तिगत लोभ अथवा चारित्रिक दुर्बलता से उपजता है।

प्रशासनिक उत्तरदायित्व के मूर्धन्य सिद्धांतकार कार्ल फ्रेडरिक \citep{friedrich1940} ने तर्क दिया था कि वृहद् प्रशासनिक संरचनाओं में लोकसेवक के आचरण का सबसे विश्वसनीय नियामक उसका 'आंतरिक नैतिक नियंत्रण' (Inner Check) और व्यावसायिक सदाचार होता है। फ्रेडरिक की मान्यता थी कि केवल बाह्य कानूनी नियम भ्रष्टाचार को रोकने में अपर्याप्त हैं; जब तक अधिकारी में अंतर्निहित नैतिक चेतना और जनसेवा की व्यक्तिगत भावना न हो, तब तक स्वच्छ प्रशासन असंभव है। इसी विमर्श को आगे बढ़ाते हुए राजनीतिक वैज्ञानिक अर्नोल्ड हाइडेनहाइमर \citep{heidenheimer1970} ने भ्रष्टाचार की शास्त्रीय परिभाषाओं को 'जनहित-केंद्रित' दृष्टिकोण के रूप में संहिताबद्ध किया, जहाँ व्यक्तिगत स्वार्थ के वशीभूत होकर किए गए निर्णयों को प्राथमिक व्याधि माना गया।

स्वतंत्र भारत के नीति-निर्माताओं ने भी आरंभिक वर्षों में इसी चारित्रिक-नैतिक दृष्टिकोण को नीति का आधार बनाया। स्वतंत्रता प्राप्ति के तत्काल पश्चात भारत सरकार द्वारा नियुक्त ए.डी. गोरवाला समिति \citep{gorwala1951} ने अपने आधिकारिक प्रतिवेदन में स्पष्ट किया कि प्रशासन की सबसे गंभीर चुनौती लोकसेवकों के नैतिक चरित्र का पतन है। समिति का मानना था कि यदि लोकसेवकों में आध्यात्मिक व नैतिक मूल्यों का पुनरुत्थान किया जाए तथा सत्यनिष्ठा प्रतिज्ञाएं लागू की जाएं, तो प्रशासनिक तंत्र को पारदर्शी बनाया जा सकता है। इसी वैचारिक दर्शन के तहत आरंभिक दशकों में आचार संहिताओं और नैतिक प्रशिक्षण पर भारी प्रशासनिक भरोसा व्यक्त किया गया।

\subsection{नैतिक दृष्टिकोण की अनुभवजन्य सीमाएँ एवं संरचनात्मक आवश्यकता}

यद्यपि व्यक्तिगत सदाचार का सैद्धांतिक महत्व निर्विवाद है, आधुनिक संस्थागत अर्थशास्त्र और लोक चयन सिद्धांत (Public Choice Theory) यह प्रमाणित करते हैं कि केवल नैतिक उपदेश प्रशासनिक भ्रष्टाचार को रोकने में अनुभवजन्य रूप से अपर्याप्त रहे हैं। \citep{shleifer1993} ने स्पष्ट किया कि जब प्रशासनिक तंत्र में अनुमतियों और लाइसेंसों का वितरण सरकारी अधिकारियों के एकाधिकार में होता है, तो भ्रष्टाचार व्यक्तिगत चरित्र का दोष न होकर बाज़ार शक्तियों और आर्थिक प्रोत्साहनों का स्वाभाविक परिणाम बन जाता है। 

इसी प्रकार, \citep{bardhan1997} और \citep{rothstein2011} के अध्ययनों ने स्थापित किया कि जहाँ पूरी प्रशासनिक व्यवस्था में अविश्वास और सूचना की विषमता (Information Asymmetry) व्याप्त हो, वहाँ भ्रष्टाचार व्यक्तिगत नैतिकता के अभाव के बजाय एक 'सामूहिक कार्रवाई समस्या' (Collective Action Problem) का रूप ले लेता है। ऐसी व्यवस्था में ईमानदारी बनाए रखने वाले व्यक्ति को संस्थागत रूप से आर्थिक और समय की भारी क्षति उठानी पड़ती है। अतः व्यक्तिगत चरित्र पर अत्यधिक बल देने के स्थान पर यह आवश्यक हो जाता है कि हम उन अंतर्निहित आर्थिक नियमों और रणनीतिक अंतःक्रियाओं का विश्लेषण करें जो सामान्य व्यक्तियों को भ्रष्ट आचरण की ओर प्रवृत्त करते हैं।

\subsection{किराया-वसूली (Rent-Seeking) एवं उपयोगिता फलन का गणितीय ढाँचा}

\subsubsection{किराया-वसूली (Rent-Seeking) की सैद्धांतिक व्याख्या}
सामान्य आर्थिक व्यवहार में धन का सृजन किसी नवीन वस्तु, सेवा अथवा मूल्य के उत्पादन से होता है। इसके विपरीत, अर्थशास्त्र में 'किराया-वसूली' उस स्थिति को कहा जाता है जहाँ कोई व्यक्ति या अधिकारी समाज में किसी नवीन मूल्य का निर्माण किए बिना, केवल अपने पद, कानूनी शक्ति अथवा प्रशासनिक एकाधिकार का उपयोग करके नागरिकों से अनुचित धन वसूलता है। 

अर्थशास्त्री रॉबर्ट क्लिटगार्ड \citep{klitgaard1988} ने भ्रष्टाचार के इस संरचनात्मक उद्भव को एक बीजगणितीय समीकरण द्वारा प्रस्तुत किया:
\begin{equation}
\text{भ्रष्टाचार } (C) = \text{एकाधिकार } (M) + \text{विवेकाधिकार } (D) - \text{जवाबदेही } (A)
\end{equation}
यहाँ $M$ (Monopoly) का अर्थ है कि किसी विशेष नागरिक सेवा (जैसे लाइसेंस, भूमि नामांतरण या जाति प्रमाण पत्र) को प्रदान करने का सरकारी तंत्र के पास पूर्ण एकाधिकार है; $D$ (Discretion) का अर्थ है कि अधिकारी के पास फाइल को मनमाने ढंग से रोकने या स्वीकृत करने का विवेकाधिकार है; तथा $A$ (Accountability) बाह्य पारदर्शी निगरानी तंत्र को दर्शाता है। क्लिटगार्ड का यह सूत्र स्पष्ट करता है कि जब तक किसी पद पर विवेकाधिकार ($D$) असीमित रहेगा और जवाबदेही ($A$) क्षीण रहेगी, तब तक केवल नैतिक व्याख्यान भ्रष्टाचार ($C$) को समाप्त नहीं कर सकते।

\subsubsection{गैरी बेकर का उपयोगिता फलन (Utility Function)}
अपराध के आर्थिक विश्लेषण में नोबेल पुरस्कार विजेता गैरी बेकर \citep{becker1968} ने प्रतिपादित किया कि कोई भी व्यक्ति अपराध या बेईमानी का चयन तभी करता है जब उससे मिलने वाला प्रत्याशित लाभ पकड़े जाने के संभावित नुकसान से अधिक हो। लोकसेवक के इस निर्णय को प्रत्याशित उपयोगिता फलन ($\mathbb{E}[U]$) द्वारा मॉडल किया जा सकता है:
\begin{equation}
\mathbb{E}[U] = (1 - P_d) \cdot U(W + B) + P_d \cdot U(W + B - F)
\end{equation}
जहाँ:
\begin{itemize}
    \item $W$ अधिकारी का निश्चित कानूनी वेतन है,
    \item $B$ अवैध परितोषण अथवा रिश्वत की राशि है,
    \item $P_d$ अधिकारी के पकड़े जाने, सिद्धदोष होने और दंडित होने की संभावना है ($0 \le P_d \le 1$),
    \item $F$ पकड़े जाने की स्थिति में होने वाला वित्तीय दंड, सेवा समाप्ति तथा सामाजिक कलंक की संयुक्त क्षति है।
\end{itemize}

जब अदालती मुकदमों में दशकों का विलंब हो और जांच एजेंसियों की पहुंच सीमित हो, तो पकड़े जाने और दंडित होने की संभावना शून्य के करीब पहुँच जाती है ($P_d \to 0$)। ऐसी परिस्थिति में $U(W + B) > U(W)$ होने के कारण अधिकारी के लिए बेईमानी करना एक आर्थिक रूप से तर्कसंगत (Rational) चयन बन जाता है।

\subsection{खेल सिद्धांत (Game Theory) एवं उप-इष्टतम नैश संतुलन}

\subsubsection{खेल सिद्धांत का सामान्य परिचय}
खेल सिद्धांत (Game Theory) गणित और अर्थशास्त्र की वह शाखा है जो यह अध्ययन करती है कि जब दो या दो से अधिक पक्ष किसी परस्पर-निर्भर स्थिति में आमने-सामने हों, तो वे एक-दूसरे के संभावित कदमों का अनुमान लगाकर अपने लाभ को अधिकतम करने के लिए क्या निर्णय लेते हैं। 

इसे सड़क पर यातायात के उदाहरण से समझा जा सकता है: यदि किसी चौराहे पर सभी वाहन चालक लाल बत्ती तोड़ रहे हों, तो अकेला चालक जो बत्ती पर रुकता है, न केवल अपनी मंजिल पर सबसे देर से पहुँचेगा, बल्कि पीछे से दुर्घटनाग्रस्त होने का जोखिम भी उठाएगा। परिणामस्वरूप, दूसरों के व्यवहार को देखकर वह भी बत्ती तोड़ने पर विवश हो जाता है। प्रशासनिक कार्यालय में नागरिक और लोकसेवक के मध्य चलने वाली अंतःक्रिया भी ठीक इसी सिद्धांत पर कार्य करती है।

\subsubsection{रणनीतिक पे-ऑफ मैट्रिक्स एवं नैश संतुलन}
प्रशासनिक पटल पर होने वाली इस अंतःक्रिया को एक $2 \times 2$ नॉर्मल-फॉर्म रणनीतिक खेल के रूप में औपचारिक रूप दिया जा सकता है। 

मान लीजिए कि नागरिक के पास दो रणनीतियाँ हैं: $\{ \text{नियम पालन } (C_1), \text{रिश्वत भुगतान } (C_2) \}$ और अधिकारी के पास दो रणनीतियाँ हैं: $\{ \text{कर्तव्यनिष्ठ निष्पादन } (O_1), \text{किराया वसूली / विलंब } (O_2) \}$। चरों का मानकीकरण निम्नलिखित है:
\begin{itemize}
    \item $V$: नागरिक के लिए उस कार्य अथवा प्रमाण पत्र का आर्थिक/सामाजिक मूल्य।
    \item $W$: अधिकारी का निश्चित कानूनी मासिक वेतन।
    \item $B$: मांगी जाने वाली रिश्वत की राशि।
    \item $t$: कार्य निष्पादन की सामान्य प्रक्रियात्मक लागत।
    \item $C_d$: अधिकारी द्वारा जानबूझकर किए गए प्रक्रियात्मक विलंब से नागरिक को होने वाली आर्थिक व मानसिक क्षति ($C_d \gg B$)।
\end{itemize}

\vspace{0.5\baselineskip}
\begin{table}[htbp]
\centering
\caption{नागरिक एवं लोकसेवक के मध्य रणनीतिक पे-ऑफ मैट्रिक्स (Game Payoff Matrix)}
\label{tab:game_payoff}
\begin{tabular*}{\textwidth}{@{\extracolsep{\fill}}lcc}
\toprule
\textbf{अधिकारी $\backslash$ नागरिक} & \textbf{नागरिक: नियम पालन ($C_1$)} & \textbf{नागरिक: रिश्वत भुगतान ($C_2$)} \\
\midrule
\textbf{अधिकारी: कर्तव्यनिष्ठ ($O_1$)} & $(V - t, \; W)$ & $(V - t - B, \; W)$ \\
\addlinespace
\textbf{अधिकारी: किराया वसूली ($O_2$)} & $(-C_d, \; W)$ & $(V - t - B, \; W + B)$ \\
\bottomrule
\end{tabular*}
\vspace{0.2em}
\begin{flushleft}
\footnotesize \textit{नोट (Note).} प्रत्येक कोष्ठक में प्रथम पद नागरिक का उपयोगिता पे-ऑफ ($U_c$) और द्वितीय पद अधिकारी का पे-ऑफ ($U_o$) दर्शाता है।
\end{flushleft}
\end{table}

\noindent \textbf{रणनीतिक संतुलन का विश्लेषण:} 
जॉन नैश द्वारा प्रतिपादित 'नैश संतुलन' (Nash Equilibrium) उस स्थिति को कहते हैं जहाँ खेल में शामिल कोई भी खिलाड़ी दूसरे खिलाड़ी की रणनीति को जानते हुए, अपने निर्णय को अकेले बदलकर अपनी स्थिति को बेहतर नहीं कर सकता। 

यदि अधिकारी किराया-वसूली रणनीति ($O_2$) का चयन कर फाइल को अटका देता है, तो नागरिक के सम्मुख दो विकल्प शेष रहते हैं: नियम पालन पर अड़े रहकर अंतहीन चक्कर काटना (जिसमें उसे $-C_d$ की भारी क्षति होती है), अथवा $B$ राशि देकर अपना कार्य समय पर करवा लेना ($V - t - B$ उपयोगिता)। चूँकि कार्यालय के चक्कर काटने की लागत रिश्वत की राशि से कई गुना अधिक भारी पड़ती है ($(V - t - B) > -C_d$), इसलिए नागरिक के लिए विवश होकर रिश्वत देना ही सर्वश्रेष्ठ रणनीतिक प्रतिक्रिया (Best Response) बन जाता है। 

दूसरी ओर, अधिकारी के लिए $(W + B) > W$ होने के कारण घूस लेना प्रभावी रणनीति बन जाता है। इस प्रकार पूरी व्यवस्था $(C_2, O_2)$ पर आकर जम जाती है। यह एक \textbf{उप-इष्टतम नैश संतुलन (Sub-optimal Nash Equilibrium)} है, जहाँ दोनों पक्ष यह जानते हैं कि रिश्वतखोरी नैतिक रूप से गलत और संपूर्ण समाज के लिए हानिकारक है, फिर भी वे तात्कालिक लाभ-हानि के चक्रव्यूह में फँसकर वही करने पर विवश हो जाते हैं \citep{aidt2003, bardhan1997, mishra2006}।

\subsection{संज्ञानात्मक पूर्वाग्रह एवं समाजशास्त्रीय आयाम}

भ्रष्टाचार केवल नकद रुपयों का आर्थिक खेल नहीं है; इसके पीछे मानवीय संज्ञान और सामाजिक व्यवस्था के गहरे मनोवैज्ञानिक तंत्र कार्य करते हैं:

\subsubsection{डूबी हुई लागत का भ्रम (Sunk Cost Fallacy)} 
निर्णय-विज्ञान में \citep{arkes1985} ने प्रतिपादित किया कि जब कोई व्यक्ति किसी पद या लक्ष्य को पाने के लिए अपना अत्यधिक समय, मानसिक ऊर्जा और वित्तीय संसाधन निवेश कर देता है, तो उसका मस्तिष्क उस खर्च हो चुके निवेश (Sunk Cost) को किसी भी तरह 'वसूलने' का प्रयास करता है। भारत में सिविल सेवा जैसी प्रतियोगी परीक्षाओं (UPSC व राज्य लोक सेवा आयोग) में चयन अनुपात $०.१\%$ से भी कम है। अभ्यर्थी अपने जीवन के स्वर्णिम ५ से ७ वर्ष एकाकी कमरों में अत्यधिक मानसिक तनाव में व्यतीत करते हैं। पद पर पहुँचने के बाद, यह पूर्व निवेश एक अवचेतन मानसिक ग्रंथि में बदल जाता है जिसे मनोविज्ञान में 'क्षतिपूर्ति का अधिकार' \citep[Psychological Entitlement;][]{campbell2004} कहा जाता है। अधिकारी के अवचेतन में यह भाव घर कर जाता है: ``मैंने इस कुर्सी को पाने हेतु अपनी युवावस्था की आहुति दी है, अतः व्यवस्था से अतिरिक्त व्यक्तिगत लाभ वसूलना मेरा स्वाभाविक अधिकार है।''

\subsubsection{नैतिक अलगाव (Moral Disengagement)} 
सत्ता और अधिकार प्राप्त होते ही व्यक्ति का दूसरों के प्रति संवेदनशीलता स्तर कम होने लगता है \citep{keltner2003}। विख्यात मनोवैज्ञानिक अल्बर्ट बंडूरा \citep{bandura1999} का सिद्धांत दर्शाता है कि कोई भी लोकसेवक स्वयं को भीतर से बुरा इंसान नहीं मानता। अपने अपराध-बोध से बचने के लिए वह चार मानसिक रक्षा-युक्तियों (Coping Mechanisms) का सहारा लेता है:
\begin{itemize}
    \item \textbf{नैतिक औचित्य (Moral Justification):} ``यह पैसा मैं अपने बच्चों के भविष्य और सामाजिक दायित्वों के निर्वहन हेतु ले रहा हूँ।''
    \item \textbf{उत्तरदायित्व का विस्थापन (Displacement of Responsibility):} ``जब ऊपर से नीचे तक पूरी व्यवस्था ही ऐसी है, तो मेरे अकेले साधु बनने से क्या व्यवस्था बदल जाएगी?''
    \item \textbf{हानि की उपेक्षा (Disregard of Harm):} ``किसी बड़े उद्योगपति या ठेकेदार से कमीशन लेने से किसी निर्धन नागरिक का प्रत्यक्ष नुकसान तो नहीं हो रहा।''
    \item \textbf{सीमित नैतिक दृष्टि (Bounded Ethicality):} व्यक्ति अपने भ्रष्ट आचरण को सामान्य दिनचर्या का हिस्सा मानकर अपनी नैतिक पहचान को खंडित होने से बचा लेता है \citep{chugh2005}।
\end{itemize}

\subsubsection{सामंती चेतना एवं 'माई-बाप' संस्कृति} 
मैक्स वेबर \citep{weber1978} के अनुसार, एक आधुनिक नौकरशाही में पद और व्यक्ति का पृथक्करण होना चाहिए, जहाँ अधिकारी केवल नियमों का लोकसेवक বাংলাকে होता है। परंतु भारत के ऐतिहासिक प्रशासनिक संदर्भ में \citep{wade1982} का अध्ययन दर्शाता है कि औपनिवेशिक काल की सामंती 'माई-बाप' संस्कृति आज भी जीवित है। नागरिक अधिकारी को लोकसेवक न मानकर शक्ति का स्वामी मानता है। जब समाज में ईमानदारी से सादा जीवन जीने वाले व्यक्ति को अक्षम माना जाए और अवैध धन से अर्जित वैभव व गाड़ियों को 'रुतबे और सफलता' का पैमाना मान लिया जाए, तो समाज स्वयं बेईमानी को एक सामाजिक मान्यता (Social Norm) प्रदान कर देता है \citep{bicchieri2017}।

\section{भारत में भ्रष्टाचार निवारण के ऐतिहासिक प्रयास एवं उनकी प्रभावशीलता}

स्वतंत्रता के पश्चात भारतीय राज्य ने भ्रष्टाचार को नियंत्रित करने हेतु कई वैधानिक व प्रशासनिक कदम उठाए हैं, जिनका विश्लेषणात्मक मूल्यांकन निम्नलिखित है:

\subsection{कौटिल्य का अर्थशास्त्र (प्राचीन संस्थागत दृष्टि)}
ईसा पूर्व चौथी शताब्दी में कौटिल्य ने अपने प्रसिद्ध ग्रंथ \textit{अर्थशास्त्र} में पहचान लिया था कि प्रशासनिक पदों पर पदस्थ व्यक्तियों में धन के प्रति स्वाभाविक आकर्षण होता है। उन्होंने लिखा था कि ``जैसे जीभ पर रखे शहद या विष को न चखना असंभव है, वैसे ही राजा के धन का प्रबंधन करने वाले अधिकारी द्वारा उसका उपभोग न करना कठिन है।'' कौटिल्य ने नैतिक उपदेशों के स्थान पर संस्थागत निगरानी, गुप्तचरों द्वारा गुप्त परीक्षण (Sting/Trap Operations) और नियमित लेखा-परीक्षण (Auditing) को अनिवार्य माना था।

\subsection{संथानम समिति (१९६२--१९६४) एवं केंद्रीय सतर्कता आयोग (CVC, १९६४)}
के. संथानम समिति की अनुशंसाओं पर देश में पहली बार संस्थागत सतर्कता तंत्र की नींव रखी गई \citep{santhanam1964}। समिति ने अत्यंत स्पष्ट शब्दों में स्वीकार किया था कि स्वतंत्रता के बाद आरंभ हुआ 'लाइसेंस-परमिट राज' और प्रशासनिक अधिकारियों को दिए गए असीमित विवेकाधिकार ही भ्रष्टाचार के वास्तविक जनक हैं।

\subsection{भ्रष्टाचार निवारण अधिनियम (PCA, १९८८ एवं २०१८ संशोधन)}
इस कानून ने लोकसेवकों द्वारा लिए जाने वाले अवैध परितोषण को एक गंभीर संज्ञेय अपराध घोषित किया। वर्ष २०१८ के संशोधन में धारा ८ के तहत पहली बार रिश्वत देने वाले व्यक्ति को भी आपराधिक दायरे में लाया गया, जिससे लेन-देन के दोनों सिरों पर अंकुश लगाया जा सके। तथापि, ईमानदार अधिकारियों को दुर्भावनापूर्ण मुकदमों से बचाने के उद्देश्य से जोड़ी गई धारा १७A (जिसमें जांच से पूर्व वरिष्ठ प्राधिकारी की पूर्व-स्वीकृति अनिवार्य की गई) ने व्यवहार में जांच एजेंसियों की गति को धीमा कर दिया है, जिससे यह प्रावधान कभी-कभी भ्रष्ट तत्वों के लिए एक कानूनी ढाल बन जाता है।

\subsection{सूचना का अधिकार अधिनियम (RTI, २००५)}
सूचना का अधिकार अधिनियम ने सरकारी फाइलों के बंद दरवाजों को खोलकर प्रशासनिक एकाधिकार पर सीधा प्रहार किया। \citep{peisakhin2010} द्वारा दिल्ली में किए गए एक यादृच्छिक क्षेत्रीय प्रयोग (Randomized Field Experiment) ने अनुभवजन्य रूप से सिद्ध किया कि राशन कार्ड प्राप्त करने में एक साधारण आरटीआई आवेदन उतना ही प्रभावी सिद्ध हुआ जितनी कि एक बिचौलिए को दी जाने वाली अवैध रिश्वत। आरटीआई ने नागरिक को अधिकारी के पटल पर निगरानी रखने की शक्ति ($A \uparrow$) प्रदान की।

\subsection{प्रत्यक्ष लाभ अंतरण (DBT) एवं डिजिटल बुनियादी ढाँचा (२०१४--वर्तमान)}
अर्थशास्त्री \citep{muralidharan2016} ने बिहार व झारखंड में सामाजिक सुरक्षा योजनाओं पर किए गए अपने बड़े पैमाने के अनुभवजन्य अध्ययन में प्रमाणित किया कि बिचौलियों और मानवीय संवितरण को हटाकर जब लाभार्थियों के बैंक खातों में आधार-सक्षम प्रत्यक्ष अंतरण (DBT) किया गया, तो व्यवस्थागत रिसाव (Leakage) में ४७\% की भारी कमी दर्ज की गई। इसने सिद्ध किया कि मानवीय विवेकाधिकार को तकनीक से प्रतिस्थापित करना नैतिक अपीलों से कहीं अधिक प्रभावी है।

\section{प्राथमिक अनुभवजन्य क्षेत्रीय अध्ययन (Primary Empirical Study)}

ऐतिहासिक वैधानिक पहलों और सैद्धांतिक मॉडलों का ज़मीनी स्तर पर परीक्षण करने हेतु जबलपुर (मध्य प्रदेश) के एक उच्चतर माध्यमिक विद्यालयी पारिस्थितिकी तंत्र में प्रत्यक्ष क्षेत्रीय सर्वेक्षण आयोजित किया गया।

\subsection{पद्धति एवं प्रतिदर्श संरचना}
\begin{itemize}
    \item \textbf{प्रतिदर्श आकार ($N$):} ६० उत्तरदाता।
    \item \textbf{स्तरीकृत उप-समूह (Stratified Sub-groups):} वरिष्ठ माध्यमिक छात्र ($n_1 = 25$), विद्यालयीन शिक्षक ($n_2 = 15$), तथा पालक व गैर-शैक्षणिक कर्मचारी ($n_3 = 20$)।
    \item \textbf{उपकरण (Instrument):} पाँच प्रमुख चरों पर आधारित संरचित द्विभाषी प्रश्नावली एवं प्रत्यक्ष साक्षात्कार।
\end{itemize}

\subsection{सांख्यिकीय परिणाम एवं क्रॉस-टैबुलेशन}

\begin{table}[htbp]
\centering
\caption{प्राथमिक क्षेत्रीय सर्वेक्षण: समग्र एवं उप-समूहवार परिणाम ($N = 60$)}
\label{tab:survey_results}
\begin{tabular*}{\textwidth}{@{\extracolsep{\fill}}lcccc}
\toprule
\textbf{सर्वेक्षण चर / प्रश्न} & \textbf{छात्र ($n=25$)} & \textbf{शिक्षक ($n=15$)} & \textbf{पालक ($n=20$)} & \textbf{कुल योग (\%)} \\
\midrule
१. प्रशासनिक कार्य में असामान्य परेशानी/विलंब & १६ ($६४.०\%$) & १० ($६६.७\%$) & १६ ($८०.०\%$) & \textbf{४२ ($७०.०\%$)} \\
२. बिना तार्किक कारण अनुचित पक्षपात/प्राथमिकता & १५ ($६०.०\%$) & ११ ($७३.३\%$) & १३ ($६५.०\%$) & \textbf{३९ ($६५.०\%$)} \\
३. अपरोक्ष घूसखोरी हेतु प्रक्रियात्मक बाध्यता & १४ ($५६.०\%$) & १० ($६६.७\%$) & १३ ($६५.०\%$) & \textbf{३७ ($६१.७\%$)} \\
\bottomrule
\end{tabular*}
\vspace{0.2em}
\begin{flushleft}
\footnotesize \textit{नोट (Note).} प्राथमिक क्षेत्रीय सर्वेक्षण आँकड़े, प्रतिभा सम्मान योजना (PSY) शोध परियोजना के अंतर्गत संकलित।
\end{flushleft}
\end{table}

\subsection{खुले प्रश्नों का गुणात्मक एवं विषयगत विश्लेषण}
\begin{itemize}
    \item \textbf{बाध्यता बनाम स्वेच्छा (प्रश्न ४):} सर्वेक्षण में ६३.३\% ($n=38$) उत्तरदाताओं ने दृढ़ता से कहा कि एक आम नागरिक घूस देने के लिए 'बाध्य' होता है, क्योंकि दफ्तरों के चक्कर काटने से उसकी दैनिक आजीविका और समय का जो नुकसान होता है, वह रिश्वत की राशि से अधिक भारी पड़ता है (यह सीधे तालिका \ref{tab:game_payoff} के खेल-सैद्धांतिक मॉडल की पुष्टि करता है)। दूसरी ओर, ३६.७\% ($n=22$) उत्तरदाताओं---विशेषकर संपन्न पालकों व व्यापारियों---ने स्वीकार किया कि कई लोग कतार तोड़ने और अवैध लाभ पाने हेतु 'स्वेच्छा' से रिश्वत की पेशकश करते हैं।
    \item \textbf{सुझाए गए समाधान (प्रश्न ५):} ८५\% उत्तरदाताओं ने माना कि जब तक सरकारी दफ्तरों में नागरिकों का अधिकारियों से 'सीधा मानवीय संपर्क' रहेगा, तब तक घूसखोरी नहीं रुकेगी; अतः फेसलेस डिजिटल व्यवस्था और शिकायतकर्ता की १००\% गोपनीय सुरक्षा ही सर्वोत्तम उपाय हैं।
\end{itemize}

\section{शोध की अंतर्निहित सीमाएँ (Limitations)}

अकादमिक सत्यनिष्ठा और निष्पक्षता के अनुपालन में अध्ययन की निम्नलिखित सीमाएं स्वीकार की जाती हैं:
\begin{itemize}
    \item \textbf{स्थानीयकृत प्रतिदर्श:} सर्वेक्षण $N=60$ के स्थानीय शैक्षणिक आकार तक सीमित है; अतः इसे पूरे भारत के वृहद् परिदृश्य का पूर्ण सामान्यीकृत निष्कर्ष नहीं माना जा सकता।
    \item \textbf{सामाजिक वांछनीयता पूर्वाग्रह (Self-Reporting Bias):} \citep{bertrand2001} के अनुसार, भ्रष्टाचार जैसे संवेदनशील और कानूनी अपराध से जुड़े विषय पर लोग कानूनी भय अथवा संकोच के कारण अपने वास्तविक आचरण को पूर्णतः साझा नहीं करते।
    \item \textbf{अपराध के गुप्त आँकड़े (Dark Figure of Crime):} \citep{biderman1967} ने रेखांकित किया था कि दोनों पक्षों की मौन सहमति से बंद कमरों में होने वाले गुप्त वित्तीय लेन-देन कभी भी आधिकारिक आंकड़ों या सामान्य साक्षात्कारों में पूर्णतः दर्ज नहीं हो पाते।
\end{itemize}

\section{नीतिगत समाधान: प्रोत्साहन संरचना का पुनर्गठन}

भ्रष्टाचार निवारण का वास्तविक समाधान लोकसेवकों के लाभ-हानि समीकरण को इस प्रकार पुनर्गठित करने में है कि ईमानदारी बनाए रखना उनके लिए सबसे सुरक्षित और आर्थिक रूप से श्रेष्ठ विकल्प बन जाए।

\subsection{प्रशासनिक विवेकाधिकार का विलोपन एवं स्वतः-स्वीकृति प्रणाली}
\begin{itemize}
    \item \textbf{एल्गोरिद्मिक फाइल आवंटन (FIFO):} सरकारी कार्यालयों में फाइलों का वितरण 'फर्स्ट-इन-फर्स्ट-आउट' (First-In-First-Out) के आधार पर स्वचालित डिजिटल सॉफ्टवेयर द्वारा किया जाए। किसी अधिकारी के पास यह चुनने की मानवीय छूट न हो कि वह किसकी फाइल पहले देखेगा।
    \item \textbf{डीम्ड अप्रूवल (Deemed Approval Mechanism):} मध्य प्रदेश लोक सेवा गारंटी अधिनियम (२०१०) की भावना के अनुरूप, प्रत्येक नागरिक सेवा हेतु समय-सीमा (उदा. ७२ घंटे) निर्धारित हो। यदि अधिकारी उस अवधि में वैध कारण दर्ज किए बिना फाइल आगे नहीं बढ़ाता, तो वह स्वतः स्वीकृत मानी जाए और जिम्मेदार अधिकारी के मासिक वेतन से प्रतिदिन के हिसाब से कटौती की जाए।
\end{itemize}

\subsection{वित्तीय पुरस्कार एवं व्हिसलब्लोअर संरक्षण ढाँचा}
\citep{dyck2010} के शोध के अनुसार, संस्थागत भ्रष्टाचार को उजागर करने में आंतरिक मुखबिरी (Whistleblowing) सबसे प्रभावी साधन है। संयुक्त राज्य अमेरिका के \textit{False Claims Act} के सफल मॉडल की तर्ज पर, भारत में व्हिसलब्लोअर को बरामद की गई अवैध संपत्ति का १०\% से १५\% वित्तीय पुरस्कार के रूप में दिया जाए तथा उसकी पहचान को ब्लॉकचेन-आधारित गोपनीय तंत्र द्वारा सुरक्षित रखा जाए।

\subsection{प्रशासनिक प्रशिक्षण में खेल सिद्धांत एवं सामाजिक प्रतिष्ठा का क्षरण}
\begin{itemize}
    \item \textbf{LBSNAA में प्रोत्साहन-ढांचा प्रशिक्षण:} लाल बहादुर शास्त्री राष्ट्रीय प्रशासन अकादमी (मसूरी) में भावी अधिकारियों को केवल सतही नैतिकता पढ़ाने के बजाय खेल सिद्धांत सिखाया जाए। उन्हें गणितीय रूप से दर्शाया जाए कि एक भ्रष्ट निर्णय कैसे समाज के पारस्परिक विश्वास को नष्ट कर अंततः उनके अपने परिवारों के लिए भी असुरक्षित वातावरण तैयार करता है।
    \item \textbf{सामाजिक रुतबे का विखंडन:} \citep{fisman2007} के अनुसार, सामाजिक मानदंड भ्रष्ट आचरण को वैधता प्रदान करते हैं। दोषसिद्ध भ्रष्ट लोकसेवकों के नाम और संपत्तियों को सार्वजनिक डिजिटल पोर्टलों पर स्थायी रूप से प्रदर्शित किया जाए ताकि उनका अवैध सामाजिक रुतबे कलंक में परिवर्तित हो सके।
\end{itemize}

\section{निष्कर्ष}

भ्रष्टाचार निवारण केवल नैतिक नारों अथवा भावनात्मक अभियानों से संभव नहीं है। ऐतिहासिक विश्लेषण और क्षेत्रीय डेटा यह प्रमाणित करते हैं कि जब तक व्यवस्था में प्रशासनिक एकाधिकार और विवेकाधिकार विद्यमान रहेगा, तब तक नागरिक और अधिकारी दोनों के लिए बेईमानी ही सबसे तर्कसंगत रणनीति बनी रहेगी।

भारत में भ्रष्टाचार पर स्थायी नियंत्रण तभी स्थापित हो सकता है जब नीति-निर्माता क्लिटगार्ड और बेकर के समीकरणों के अनुरूप प्रोत्साहन संरचना को पुनर्गठित करें---जहाँ विवेकाधिकार शून्य हो, फेसलेस डिजिटल तकनीक द्वारा मानवीय संपर्क न्यूनतम हो, पकड़े जाने की निश्चितता ($P_d \to 1$) हो, और व्हिसलब्लोअर को संस्थागत वित्तीय सुरक्षा प्राप्त हो। केवल इसी संरचनात्मक रूपांतरण से ईमानदारी प्रशासनिक तंत्र का स्वाभाविक और गणितीय रूप से सर्वश्रेष्ठ विकल्प बन सकेगी।

\newpage

% --- REFERENCES (APA 7th Standard via Natbib / Bibitem) ---
\begin{thebibliography}{99}

\bibitem[Aidt(2003)]{aidt2003}
Aidt, T. S. (2003).
\newblock Economic analysis of corruption: A survey.
\newblock \emph{The Economic Journal}, \emph{113}(491), F632--F652.
\newblock https://doi.org/10.1111/1468-0297.00167

\bibitem[Arkes \& Blumer(1985)]{arkes1985}
Arkes, H. R., \& Blumer, C. (1985).
\newblock The psychology of sunk cost.
\newblock \emph{Organizational Behavior and Human Decision Processes}, \emph{35}(1), 124--140.
\newblock https://doi.org/10.1016/0749-5978(85)90049-4

\bibitem[Bandura(1999)]{bandura1999}
Bandura, A. (1999).
\newblock Moral disengagement in the perpetration of inhumanities.
\newblock \emph{Personality and Social Psychology Review}, \emph{3}(3), 193--209.
\newblock https://doi.org/10.1207/s15327957pspr0303\_3

\bibitem[Bardhan(1997)]{bardhan1997}
Bardhan, P. (1997).
\newblock Corruption and development: A review of issues.
\newblock \emph{Journal of Economic Literature}, \emph{35}(3), 1320--1346.

\bibitem[Becker(1968)]{becker1968}
Becker, G. S. (1968).
\newblock Crime and punishment: An economic approach.
\newblock \emph{Journal of Political Economy}, \emph{76}(2), 169--217.
\newblock https://doi.org/10.1086/259394

\bibitem[Bertrand \& Mullainathan(2001)]{bertrand2001}
Bertrand, M., \& Mullainathan, S. (2001).
\newblock Do people mean what they say? Implications for subjective survey data.
\newblock \emph{American Economic Review}, \emph{91}(2), 67--72.
\newblock https://doi.org/10.1257/aer.91.2.67

\bibitem[Bicchieri \& Ganegoda(2017)]{bicchieri2017}
Bicchieri, C., \& Ganegoda, D. B. (2017).
\newblock Determinants of corruption: A social-norms approach.
\newblock \emph{Social Research: An International Quarterly}, \emph{84}(3), 653--668.

\bibitem[Biderman \& Reiss(1967)]{biderman1967}
Biderman, A. D., \& Reiss, A. J. (1967).
\newblock On exploring the ``dark figure'' of crime.
\newblock \emph{The Annals of the American Academy of Political and Social Science}, \emph{374}(1), 1--15.
\newblock https://doi.org/10.1177/000271626737400102

\bibitem[Campbell et~al.(2004)Campbell, Bonacci, Shelton, Exline, \& Bushman]{campbell2004}
Campbell, W. K., Bonacci, A. M., Shelton, J., Exline, J. J., \& Bushman, B. J. (2004).
\newblock Psychological entitlement: Interpersonal consequences and validation of a self-report measure.
\newblock \emph{Journal of Personality Assessment}, \emph{83}(1), 29--45.
\newblock https://doi.org/10.1207/s15327752jpa8301\_04

\bibitem[Chugh et~al.(2005)Chugh, Bazerman, \& Banaji]{chugh2005}
Chugh, D., Bazerman, M. H., \& Banaji, M. R. (2005).
\newblock Bounded ethicality as a psychological barrier to recognizing conflicts of interest.
\newblock In D. A. Moore, D. M. Cain, G. Loewenstein, \& M. H. Bazerman (Eds.), \emph{Conflicts of interest: Challenges and solutions in business, law, medicine, and public policy} (pp. 74--95). Cambridge University Press.

\bibitem[Dobel(1978)]{dobel1978}
Dobel, J. P. (1978).
\newblock The corruption of a state.
\newblock \emph{American Political Science Review}, \emph{72}(3), 958--973.
\newblock https://doi.org/10.2307/1955113

\bibitem[Dyck et~al.(2010)Dyck, Morse, \& Zingales]{dyck2010}
Dyck, A., Morse, A., \& Zingales, L. (2010).
\newblock Who blows the whistle on corporate fraud?
\newblock \emph{The Journal of Finance}, \emph{65}(6), 2213--2253.
\newblock https://doi.org/10.1111/j.1540-6261.2010.01614.x

\bibitem[Fisman \& Miguel(2007)]{fisman2007}
Fisman, R., \& Miguel, E. (2007).
\newblock Corruption, norms, and legal enforcement: Evidence from diplomatic parking tickets.
\newblock \emph{Journal of Political Economy}, \emph{115}(6), 1020--1048.
\newblock https://doi.org/10.1086/527495

\bibitem[Friedrich(1940)]{friedrich1940}
Friedrich, C. J. (1940).
\newblock Public policy and the nature of administrative responsibility.
\newblock \emph{Public Policy}, \emph{1}, 3--24.

\bibitem[Gorwala(1951)]{gorwala1951}
Gorwala, A. D. (1951).
\newblock \emph{Report on public administration}.
\newblock Planning Commission, Government of India.

\bibitem[Heidenheimer(1970)]{heidenheimer1970}
Heidenheimer, A. J. (Ed.). (1970).
\newblock \emph{Political corruption: Readings in comparative analysis}.
\newblock Holt, Rinehart and Winston.

\bibitem[Keltner et~al.(2003)Keltner, Gruenfeld, \& Anderson]{keltner2003}
Keltner, D., Gruenfeld, D. H., \& Anderson, C. (2003).
\newblock Power, approach, and inhibition.
\newblock \emph{Psychological Review}, \emph{110}(2), 265--284.
\newblock https://doi.org/10.1037/0033-295X.110.2.265

\bibitem[Klitgaard(1988)]{klitgaard1988}
Klitgaard, R. (1988).
\newblock \emph{Controlling corruption}.
\newblock University of California Press.

\bibitem[Mishra(2006)]{mishra2006}
Mishra, A. (2006).
\newblock Persistence of corruption: Some issues.
\newblock \emph{Journal of Development Economics}, \emph{79}(2), 396--407.
\newblock https://doi.org/10.1016/j.jdeveco.2005.03.003

\bibitem[Muralidharan et~al.(2016)Muralidharan, Niehaus, \& Sukhtankar]{muralidharan2016}
Muralidharan, K., Niehaus, P., \& Sukhtankar, S. (2016).
\newblock Building state capacity: What we can learn from Bihar's social security pensions.
\newblock \emph{American Economic Review}, \emph{106}(10), 2895--2929.
\newblock https://doi.org/10.1257/aer.20141177

\bibitem[Peisakhin \& Pinto(2010)]{peisakhin2010}
Peisakhin, L., \& Pinto, P. (2010).
\newblock Is transparency an effective anti-corruption strategy? Evidence from a field experiment in India.
\newblock \emph{Regulation \& Governance}, \emph{4}(3), 261--280.
\newblock https://doi.org/10.1111/j.1748-5991.2010.01081.x

\bibitem[Philp(1997)]{philp1997}
Philp, M. (1997).
\newblock Defining political corruption.
\newblock \emph{Political Studies}, \emph{45}(3), 436--462.
\newblock https://doi.org/10.1111/1467-9248.00090

\bibitem[Rothstein(2011)]{rothstein2011}
Rothstein, B. (2011).
\newblock \emph{The quality of government: Corruption, social trust, and inequality in international perspective}.
\newblock University of Chicago Press.

\bibitem[Santhanam(1964)]{santhanam1964}
Santhanam, K. (1964).
\newblock \emph{Report of the Committee on Prevention of Corruption}.
\newblock Ministry of Home Affairs, Government of India.

\bibitem[Shleifer \& Vishny(1993)]{shleifer1993}
Shleifer, A., \& Vishny, R. W. (1993).
\newblock Corruption.
\newblock \emph{The Quarterly Journal of Economics}, \emph{108}(3), 599--617.
\newblock https://doi.org/10.2307/2118402

\bibitem[Thompson(1993)]{thompson1993}
Thompson, D. F. (1993).
\newblock Mediated corruption: The case of the Keating Five.
\newblock \emph{American Political Science Review}, \emph{87}(2), 369--381.
\newblock https://doi.org/10.2307/2939047

\bibitem[Wade(1982)]{wade1982}
Wade, R. (1982).
\newblock The system of administrative and political corruption: Canal irrigation in South India.
\newblock \emph{The Journal of Development Studies}, \emph{18}(3), 287--328.
\newblock https://doi.org/10.1080/00220388208421833

\bibitem[Weber(1978)]{weber1978}
Weber, M. (1978).
\newblock \emph{Economy and society: An outline of interpretive sociology} (G. Roth \& C. Wittich, Eds.).
\newblock University of California Press. (मूल कार्य १९२२ में प्रकाशित).

\end{thebibliography}

\end{document}

